\part{V: Identity, Change, and Worldhood}\label{part-v-identity-change-and-worldhood}

\emph{A numbering note before this part begins: the original planning
assigned Chapter 19 to Part VI. Resolving that now rather than letting it
accumulate --- Part V takes Chapters 16--19, and Part VI shifts to
Chapters 20--24. The renumbering of Part VI's own files is deferred to
that part's final revision pass, per the drafting order already settled
on; this note exists so the collision doesn't get forgotten between now
and then.}

Unlike Parts II through IV, this part is not aimed at synthesis. The
correspondence matrix already marks identity as a Bridge and flags
genuine tensions between Priest's accounts of change and motion and this
book's own machinery. Part V is where that's allowed to matter: can this
framework disagree with Priest without the disagreement collapsing back
into a redescription of him in different words? Each chapter below
presents Priest's position~\cite{priest2006} at full strength before departing from it, and
at least one departure is left as a live disagreement rather than
resolved.

\hypertarget{chapter-16-identity-under-transport}{%
\chapter*{Chapter 16 --- Identity Under Transport}
\addcontentsline{toc}{chapter}{Chapter 16 --- Identity Under Transport}
\markboth{Chapter 16 --- Identity Under Transport}{Chapter 16 --- Identity Under Transport}\label{chapter-16-identity-under-transport}}

\textbf{Priest first.} Necessary identity, contingent identity, rigid versus
non-rigid designation, and the problem of tracking a single referent
across possible worlds are old and difficult questions in their own
right, independent of anything this book adds. Rigid designation, in
particular, is a specific technical commitment: a rigid designator picks
out the same individual at every world where that individual exists,
full stop, with reference held fixed while everything else about the
world varies. Whatever this book proposes needs to earn its place beside
that account, not simply assume it's been superseded.

\textbf{The reversal, stated carefully.} Rather than beginning with two
objects \(a, b\) and asking whether \(a = b\) persists across states, begin
with a history --- an admissible continuation
\[x_{t_0} \rightsquigarrow x_{t_1} \rightsquigarrow \cdots \rightsquigarrow x_{t_n}\]
--- and ask what invariant, if any, warrants treating the endpoints as
stages of a single continuing entity rather than as a sequence of merely
resembling but numerically distinct objects. This gives a specific,
appropriately limited thesis:

\begin{quote}
Identity under transport is not primitive equality repeated across
states; it is an invariant of admissible continuation.
\end{quote}

\textbf{What this is not yet claiming.} It would overreach to say this
\emph{derives} Priest's rigid designation --- nothing here shows rigid
designation is somehow reducible to continuation-invariance, and this
chapter doesn't attempt that reduction. Rigid designation becomes,
instead, the comparison case: it holds reference fixed while worlds vary,
treating fixity as primitive; this book's account asks what structure
would make holding reference fixed a \emph{legitimate} move in the first
place, rather than an unexplained stipulation. Those are different
projects that happen to bear on the same phenomenon, and keeping them
distinct is more honest than claiming an unearned reduction.

\textbf{Reusing Chapter 2, explicitly.} Chapter 2 argued that distinction
precedes truth --- a proposition presupposes a prior act of carving out a
determinate content. The same move applies here directly: before asking
whether two appearances are stages of one identity, the system already
has to have determined what counts as a distinguishable continuation at
all, as opposed to two independent, coincidentally similar objects. This
isn't a new principle invented for this chapter; it's Chapter 2's
principle, applied one level up, from propositions to persisting things.

\textbf{An open connection to Chapter 7, surfaced rather than resolved here.}
Chapter 7 introduced \(R_w\), a regime that can vary state by state rather
than one fixed \(R\) governing everywhere. If \(R_w\) genuinely varies across
a continuation \(x_{t_0} \rightsquigarrow \cdots \rightsquigarrow x_{t_n}\),
this chapter's talk of an ``admissible continuation'' is underspecified:
admissible relative to which regime --- the source state's, the target
state's, some rule governing the transition itself? This chapter doesn't
answer that, and it shouldn't be answered by quietly picking one option
without noticing the choice was made. It's a live question for whoever
next develops this material, and it connects suggestively to Chapter
17's move of giving transitions their own type distinct from the states
they connect: if a transition step needs its own admissibility condition
rather than inheriting one from either endpoint, that would mirror
Chapter 17's ontological move rather than merely resembling it by
coincidence. Left open here, not assumed either way.

\hypertarget{chapter-17-the-instant-of-change}{%
\chapter*{Chapter 17 --- The Instant of Change}
\addcontentsline{toc}{chapter}{Chapter 17 --- The Instant of Change}
\markboth{Chapter 17 --- The Instant of Change}{Chapter 17 --- The Instant of Change}\label{chapter-17-the-instant-of-change}}

\textbf{Priest at full strength.} If something changes from \(A\) to \(\neg A\),
what characterizes the instant of transition itself? Assigning that
instant wholly to the \(A\) side, or wholly to the \(\neg A\) side, looks
arbitrary either way --- nothing distinguishes the transition instant from
any other instant on the side it's assigned to, yet something has to be
different about it, since it's specifically the boundary. Priest's
dialetheic proposal is that the transition instant genuinely instantiates
both: at that instant, and only that instant, \(A\) and \(\neg A\) both hold.
This is not a cheap move --- it's a serious answer to a genuinely hard
problem, and it deserves to be stated at this strength before anything
else is proposed.

\textbf{The rival answer.} Introduce a continuation-interface as a distinct
kind of object, rather than enriching the valuation of the boundary
state itself:
\[A \ \mid\ \partial(A, \neg A) \ \mid\ \neg A.\]
The essential move is that the boundary object \(\partial(A, \neg A)\)
need not share the same type as the states it connects. \(A\) and \(\neg A\)
are states; \(\partial(A, \neg A)\) is a transition, a different kind of
entity, not a third state squeezed between the other two and not a state
that happens to satisfy both. This is a genuinely deeper disagreement
than ``Priest says both hold; this book says there's a boundary object'' ---
Priest's account enriches the \emph{valuation} available to states (a state
can satisfy more than one of \(A\), \(\neg A\)); this book's account
proposes instead enriching the \emph{ontology of transitions} (adding a new
kind of object the formalism didn't have room for), leaving state
valuation exactly as constrained as it always was.

\begin{figure}[htbp]
\centering
\begin{tikzpicture}[
  font=\small,
  state/.style={draw, circle, minimum size=9mm, inner sep=1pt},
  iface/.style={draw, rectangle, minimum width=15mm, minimum height=9mm, inner sep=2pt},
  arr/.style={-{Stealth[length=2.2mm]}}
]
\node[state] (A1) at (0,0)    {$A$};
\node[iface] (I)  at (3,0)    {$\partial(A,\neg A)$};
\node[state] (nA1) at (6,0)   {$\neg A$};
\draw[arr] (A1) -- (I);
\draw[arr] (I) -- (nA1);
\node[anchor=west, font=\footnotesize\itshape] at (7.0,0) {typed transition};

\node[state] (A2) at (0,-2)   {$A$};
\node[state] (M)  at (3,-2)   {$A \wedge \neg A$};
\node[state] (nA2) at (6,-2)  {$\neg A$};
\draw[arr] (A2) -- (M);
\draw[arr] (M) -- (nA2);
\node[anchor=west, font=\footnotesize\itshape] at (7.0,-2) {contradictory state};
\end{tikzpicture}
\caption{Two representations of the instant of change. Above: this
book's continuation-interface, where the middle object is a
\emph{transition} --- a different type from the states it connects
(drawn as a different shape, because it is one). Below: Priest's
dialetheic account, where the middle object is a \emph{state} satisfying
both $A$ and $\neg A$. The disagreement is not over which truth value
the middle point receives; it is over whether the transition should be
represented as a state at all. Per this chapter's qualifier, neither
representation is claimed to be correct for every change.}
\label{fig:state-vs-transition}
\end{figure}

\textbf{If this holds formally, the payoff is sharp --- and needs an equally
sharp qualifier.} Should this distinction survive further scrutiny, it
suggests: a contradiction attributed to a state may sometimes indicate
that the formalism lacks an object representing the transition between
states, rather than indicating that the state itself is genuinely
overdetermined. But ``sometimes'' is doing real work in that sentence and
cannot be dropped. Part III spent four chapters establishing that
genuine gluts --- states legitimately supporting both a proposition and its
negation --- are not errors requiring elimination. This chapter's proposal
must not be read as a general strategy for explaining away every glut as
``really'' a missing transition object; that would undo Part III's central
result rather than complementing it. The honest claim is narrower: \emph{some}
contradictions attributed to states may be artifacts of a state-only
vocabulary lacking a transition object, and identifying which
contradictions are which --- genuine overdetermination versus a
missing-object artifact --- is exactly the kind of question this chapter
raises and does not settle.

\hypertarget{chapter-18-motion-without-contradictory-location}{%
\chapter*{Chapter 18 --- Motion Without Contradictory Location}
\addcontentsline{toc}{chapter}{Chapter 18 --- Motion Without Contradictory Location}
\markboth{Chapter 18 --- Motion Without Contradictory Location}{Chapter 18 --- Motion Without Contradictory Location}\label{chapter-18-motion-without-contradictory-location}}

This chapter puts RSVP to genuine argumentative work rather than treating
it as a source of illustrative vocabulary, and the correspondence stays
marked \textbf{Tension}, as the matrix already has it --- this chapter doesn't
try to resolve that rating.

\textbf{Two rival representations of motion.} Priest's Hegelian treatment
characterizes motion through contradictory occupancy: at the relevant
instant, a moving object is both at a location and not at that location,
which is how motion gets represented at all within a framework built from
static positional predicates evaluated instant by instant. RSVP offers a
different starting point entirely --- it already has a vector field \(v\),
so motion doesn't need to be reconstructed from a sequence of static
positions in the first place; it's represented dynamically from the
start:
\[(\Phi, v, S)_t \longmapsto (\Phi, v, S)_{t + \Delta t}.\]

\textbf{The question this actually raises.} Not ``which account is correct,''
but something more specific: do contradictory positional predicates
reveal something genuinely ontological about motion --- that motion really
does involve a kind of instantaneous self-contradiction --- or do they
reveal something merely \emph{inadequate} about representing motion using
static positional predicates alone, a limitation of the vocabulary rather
than a discovery about the world? RSVP's vector field doesn't answer this
question by fiat; it demonstrates that a coherent, contradiction-free
representation of motion is available once the representation is allowed
to include velocity as a primitive rather than trying to derive motion
from position sequences after the fact. Whether that shows Priest's
contradiction is dispensable, or merely shows one particular vocabulary
happens to avoid it, is the tension this chapter leaves standing.

\textbf{The backward connection to Chapter 16.} This sharpens rather than
merely echoes Chapter 16's thesis. If states are projections of
histories --- if what's available at an instant is a slice through a
continuation rather than a self-sufficient object --- then a
position-only representation at a single instant has \emph{already} discarded
exactly the information (velocity, direction, rate) needed to distinguish
``is located here'' from ``is moving through here.'' Priest's contradictory
occupancy may be less a fact about motion and more a symptom of asking a
state-only vocabulary a question it was never equipped to answer, given
that the same vocabulary was already shown, in Chapter 16, to be too thin
to carry identity across time either.

\hypertarget{chapter-19-time-memory-and-irreversibility}{%
\chapter*{Chapter 19 --- Time, Memory, and Irreversibility}
\addcontentsline{toc}{chapter}{Chapter 19 --- Time, Memory, and Irreversibility}
\markboth{Chapter 19 --- Time, Memory, and Irreversibility}{Chapter 19 --- Time, Memory, and Irreversibility}\label{chapter-19-time-memory-and-irreversibility}}

\textbf{What Chain of Memory is not being used for.} CoM formalizes reasoning
trajectories --- sequences of latent states an inferential process passes
through --- and this chapter does not derive a theory of physical time from
it. That would overreach in exactly the way this book has tried to avoid
elsewhere: CoM's trajectories are about how a reasoning process gets from
one cognitive state to another, not about the structure of time itself,
and treating the two as the same theory would be a category error, not a
correspondence.

\textbf{What CoM is used for instead: a narrower, real distinction.}
\[\text{instantaneous state} \neq \text{history-indexed state}.\]
Two states that are, at a given moment, indistinguishable in every
respect that a snapshot could capture can nonetheless have arrived by
different histories --- and those different histories can license
different continuations going forward, even though nothing about the
present moment, examined on its own, shows this. Memory, on this
account, is not simply extra information carried inside an otherwise
complete state. It can encode \emph{the path by which the state became
reachable}, which is a fact about the state's history, not a fact
addable to its instantaneous content without changing what kind of
object a ``state'' is taken to be.

\textbf{A route from Priest toward this book's irreversibility work, stated as
a route rather than an equivalence.} Priest's Spread Hypothesis and his
broader treatment of temporal flow are themselves attempts to give
content to something like ``the present moment is thicker than an
instant.'' This chapter's history-indexed-state distinction is a
different, narrower device aimed at a related target --- not offered as
the same theory in different notation, and not offered as a replacement
for Priest's account, but as a route toward this book's own
irreversibility material that happens to pass near Priest's territory
without claiming to arrive at the same destination.

\textbf{Where the Leibniz Continuity Condition tension belongs, finally.} A
discrete repair operation, \(A_t \to A_{t+1}\), can look like a
discontinuous jump at one descriptive level --- nothing between the two
states, an instantaneous replacement --- while being implemented, at a
finer level of description, by a continuous underlying trajectory that
the discrete description simply doesn't resolve. This chapter's job is to
ask, honestly, whether that distinction actually reconciles the Leibniz
Continuity Condition (change is continuous, no jumps) with this book's
discrete repair events, or whether it's merely a redescription that
doesn't touch the real disagreement. If a finer-grained continuous
trajectory really does underlie every discrete repair this book has
posited, the tension dissolves without either side having to give
ground. If it doesn't --- if some repairs are genuinely discontinuous even
under arbitrarily fine redescription --- the tension should be preserved
exactly as the correspondence matrix has kept it, rather than talked
away by an appeal to description-level relativity that hasn't actually
been shown to apply.

\hypertarget{the-arc-of-part-v-and-what-it-hands-to-part-vi}{%
\section{The arc of Part V, and what it hands to Part VI}\label{the-arc-of-part-v-and-what-it-hands-to-part-vi}}

Four chapters, one continuous progression: identity requires continuation
across states (Chapter 16); change introduces a transition object that
states alone don't contain (Chapter 17); motion requires dynamical
information absent from position predicates taken instant by instant
(Chapter 18); memory distinguishes histories that terminate in otherwise
identical states (Chapter 19). Each chapter, in turn, challenges the
sufficiency of an instantaneous state a little further, and the four
together support a thesis Part V can claim as its own:

\begin{quote}
A state does not contain all the structure required to explain its own
continuation.
\end{quote}

This is not merely a summary --- it's exactly what Part VI needs before its
machinery can be trusted. Once Part VI moves into local sections,
restriction maps, reconstruction, and gluing, there is now a stated
reason not to assume a section's pointwise values exhaust its semantics:
Part V has spent four chapters showing that assumption fails even for
single, non-localized states, well before any cover or gluing condition
enters the picture.

\textbf{Two debts to preserve explicitly for Part VI's final pass, not to
resolve here.} First, Part IV established three genuinely different
failure modes --- glut (overdetermination), incompatibility (compatibility
obstruction), and dispersion (persistence failure under coarse-graining)
--- and this book should resist deciding whether dispersion is ``really''
topology until Part VI's own machinery is on the table to test against.
If Chapter 23/24's proposed topology turns out unable to distinguish
refinement-instability from ordinary incompatibility, that's evidence the
topology needs more structure, not a license to identify the two
phenomena by definitional fiat. Second, and more general: the final Part
VI pass should function less as an occasion to add machinery and more as
a type-check across the whole book --- every occurrence of support,
construction, admissibility, accessibility, compatibility, repair,
persistence, transition, and obstruction, checked against every other
part, to make sure two notions this book worked to keep separate haven't
quietly collapsed back together somewhere along the way.

\hypertarget{status-of-this-draft}{%
\section{Status of this draft}\label{status-of-this-draft}}

Written to let disagreement stand where it's earned. Chapter 17's central
result is qualified as carefully as its payoff is stated --- ``sometimes,''
not ``always,'' and the reason that qualifier can't be dropped is spelled
out rather than left implicit. Chapter 18 keeps its Tension rating rather
than resolving it in either direction. Chapter 19 poses the LCC question
as a genuine either/or rather than assuming the reconciling answer.
Nothing in this part quietly redescribes Priest in this book's vocabulary
without either agreeing with him outright or stating a specific,
checkable point of departure --- which was the actual test this part was
supposed to be.
